\section{Background and Related Work}
This section reviews clean-label backdoor poisoning attack, sampling strategies, and defenses that are relevant to our analysis of trigger value selection.

\subsection{Clean-Label Backdoor Poisoning Attacks}
Severi et al. introduced an explanation-guided clean-label backdoor poisoning attack against malware classifiers~\cite{severi}.
Their method uses SHAP values to select influential trigger features and assigns trigger values using strategies such as MinPopulation, CountAbsSHAP, and the Combined selector.
The poisoned samples retain their benign labels, causing the trained model to associate the trigger with the benign class while maintaining normal performance on clean samples.
Other studies have examined clean-label malware backdoors under different assumptions, including black-box attacks using PE surface information and selective backdoors with realizable triggers~\cite{zheng_clean_label,jigsaw}.

\subsection{Sampling Strategies in Clean-Label Malware Poisoning}
In addition to selecting trigger features and trigger values, choosing which samples to poison can also affect the effectiveness of malware backdoor poisoning attacks.
Previous work has argued that benignware samples are not all equally useful for poisoning. Some benign samples are more similar to malware samples than others.
Selecting these malware-similar benign samples as poisoning targets can improve the effectiveness of the backdoor attack.
This approach uses similarity metrics such as feature-based distance, distribution-based distance, and contribution-space distance~\cite{pmsb}.
We use sampling strategies such as PMSB to evaluate whether the observed effectiveness--detectability relationship remains consistent under different poisoning-target selection strategies.

\subsection{Defenses and Detectability}
Previous studies have proposed defenses against clean-label backdoor poisoning attacks on malware classifiers.
Severi et al. evaluated sanitization methods such as Isolation Forest, Spectral Signature, and clustering-based detection~\cite{severi}.
Other studies have investigated data sanitization defenses based on ensemble model disagreement, nested training, dimensionality reduction, clustering, and model-agnostic mitigation strategies~\cite{ho_sanitization,venkatesan_nids,countermeasures,model_agnostic_mitigation,reddy_improved_nt,partition,yudin_retraining}.
These defenses suggest that poisoned samples can be detected through multiple signals, including data anomalies, shared directions in feature space, prediction disagreement, and clustering behavior.
In this work, we use representative sanitization defenses to measure how different value-selection strategies affect the detectability of poisoned samples.